\section{The Goals}
The main goal of coding conventions in general is to help you writing better code.

So what is \textit{good code}?

\paragraph{Good Code} does not have bugs and is running fast, without wasting valuable resources.

\bdq{Does not have bugs} also means that the code is doing what it is intended to do. We all know that it is nearly impossible to produce completely bug free code, for several reasons. One important reason is that already the specification, used as the base for the implementation, is often faulty. So the code is doing what it should do according to the specification, but it is not doing what it is meant to do~…

However, we should not forget that there is still the possibility for \textit{real} bugs.

And often you see the situation that your code is getting slower and slower, or is eating up more CPU cycles and/or memory than available. That may be due to higher data volumes or more users or lots of other reasons.

But as a result, you will have to modify your code – although, in most cases, it will be someone else who have to dig into \textit{your} code.  Or you are the unlucky one who needs to revise someone else's code.

So what you want is code that helps you with being modified – you want to have \textit{maintainable code}.

\paragraph{Maintainable Code}\index{maintainable code} supports its modification in various ways.

From a magazine dealing with software development tools, I found the following statement\footnote{I noted down the statement, but forgot to note the source – so if you know that source, I would appreciate if you would share it with me.}:

\begin{quote}
\bdq{If you compare software development with the Apollo program, most programmers would be very successful bringing man to the moon, but never will get them back alive because of their common incapability to deliver software that can be maintained over a longer period of time with reasonable costs.} \\
\textit{Unknown Source}\autocite{UnknownSource1}
\end{quote}

Maintainability has a lot of facets, and I will discuss several of them in the run of this document. But before you can even think of modifying the code – either for fixing a bug, optimising the memory consumption, or adding new functionality – you need to identify the location.  

For that you have to read the already existing code, and of course, you have to understand what you are reading. So a crucial quality of maintainable code is that it is easy to read and \textit{basically} understandable\footnote{Ok, the requirement should be \bsq{easy to understand}, but that is a relative term. If something is \bsq{easily} understandable depends on the experience level of the reader, and here we talk not only about their coding experience, but also about their familiarity with the problem domain. And you cannot always code for the rookies.}.

No matter whether your code will meet all the other criteria that would make it maintainable: if it is not readable, it is not really maintainable. 

\paragraph{Readable Code}\index{readable code} can be easily consumed by a human.

\begin{quote}
\bdq{[…], programs must be written for people to read, and only incidentally for machines to execute.}\\
Herold Abelson et. al.: \textit{Structure and Interpretation of Computer Programs}
\autocite{Sussman:StructureAndInterpretationOfComputerPrograms}
\end{quote}

\begin{quote}
\bdq{Everybody can write code that can be read by a computer, but only good developers will write code that can be read by humans.}\footnote{I agree in general, but also the converse is true: only developers that will write code that can be read by humans are good developers.} \\
Martin G. Fowler: \textit{Refactoring: Improving the Design of Existing Code}\autocite{Fowler:Refactoring}
\end{quote}

Reading code is more important than writing code.\autocite{Marks:LocalVariableTypeInference:ReadingIsMoreImportantThanWriting}

Code is read much more often than it is written. Shorter programs can be preferable to longer ones, but shortening a program too much can omit information that is useful for understanding the code. The central issue here is to find the right size for the program such that understandability is maximized.

You should be specifically unconcerned here with the amount of keyboarding that is necessary to input or to edit a program. While concision may be a nice bonus for the author, focusing on it misses the main goal, which is to improve the understandability of the resulting code.

And from my experience I found that verbosity is your friend!\footnote{Refer to the “Coding Guardrails”, bullet point \ref{lst:ZoP:ExplicitVsImplicit}.}

Several guidelines and recommendations in this document, in particular regarding how to apply comments and how to write comments, but also those about to write a method, do require significant additional typing. Java in general has the reputation to be too verbose, and this coding conventions will even add to that verbosity – but as said: that's not generally bad!

But if you are afraid of the keyboarding, you should learn typewriting! For my understanding, someone who does not reach at least 100~CPM\footnote{CPM = “Characters per minute”, or, in German: „Anschläge pro Minute“} should look for a job outside of software development!\footnote{Just for your comparison: average values for a typist or secretary are around 200 CPM, peak values are at 500 to 900 CPM (the world record is 955~CPM). See also \autocite{RATATYPE}.}

Writing code is not the only area where you would benefit from mastering that important skill; it will also help you to write all the other stuff you have to deliver in addition to code (documentation, meeting notes, emails, specification documents,~…) – from a site that offers typing courses\autocite{RATATYPE}:
\begin{displayquote}
\textbf{WHEN TO START LEARNING}

Developed typing skills can help young people get better study results in school or college, and get better job offers once they have finished school. 65\% of people who learn how to touch type faster are under the age of 24. The main goal of learning to touch type or improving touch typing skills, after age 25 is to become more successful in current job.

\textbf{MOTIVATION}

Improving touch typing skills is a must for every person under 18, but everyone can benefit from some time spent practicing. Average typing speed is one of the key skills listed on your resume. Spend some more time practicing and get better results […].
\end{displayquote}

That they do marketing to sell their product does not make their statements wrong.

Also, modern IDEs do a lot of the typing for you, as long as you take the trouble to configure them appropriately. Refer to chapter \tqfullref{sec:ToolSupport}.

Code readability shouldn’t depend on IDEs.\autocite{Marks:LocalVariableTypeInference:NoIDEDependency}

Code is often written and read within an IDE, so it is tempting to rely heavily on the code analysis features of IDEs. Unfortunately, code is at least as often read \textit{outside} an IDE. It appears in many places where IDE facilities are not available, such as snippets within a document, browsing a repository on the internet, or in a patch file. It is counterproductive to have to import code into an IDE simply to understand what the code does.

Code should be self-revealing. It should be understandable on its face, without the need for assistance from tools.

\subsection{Conclusion}
To summarize from above, the goal for these coding conventions is to help you to write~…
\begin{itemize}[nosep]
	\item …~good,~…
	\item …~maintainable,~…
	\item …~readable~…
\end{itemize}
…~Java source code.

As a side effect, all source code looks familiar to all members of the team. That conformity further supports readability, but such a common look also identifies the source code as written by the team, as its trademark.

%==============================================================================
% End of file!!
%==============================================================================
